\begin{table*}[htbp]
\centering\small
\caption{Categorias temáticas selecionadas.}
\label{tab:usuarios-temas}
\begin{tabular}{p{0.20\textwidth}cp{0.30\textwidth}p{0.34\textwidth}}
\hline
Categoria & n & Trecho literal sem identificação & Implicação\\ \hline
Representação visual e variabilidade espacial & 16 & ``Ver o mapa espacial e o gráfico temporal na mesma tela'' & Manter integração visual e comunicar limites dos índices e das cores.\\
Navegação, desenho e orientação & 17 & ``Não consegui desenhar a área sozinho'' & Priorizar orientação inicial e revisão das tarefas de demarcação e exportação.\\
Transparência e parametrização dos alertas & 13 & ``Falta clareza sobre a regra estatística'' & Explicar janela e limiar; não transformar sinalização em diagnóstico.\\
Qualidade e proveniência dos dados & 11 & ``Informar a fonte e a resolução dos dados meteorológicos'' & Expor metadados; esclarecer que NASA POWER não é estação local.\\
Latência e conectividade percebidas & 3 & ``Lentidão no carregamento de camadas raster'' & Medir latência e transferência antes de atribuir causas ou ganhos.\\
Uso móvel e acesso offline & 3 & ``Desenvolver versão otimizada para celular'' & Testar dispositivos físicos; estudar offline sem confundir com função atual.\\
\hline
\end{tabular}
\fonteTCC{dados da pesquisa (2026).}
\notaTCC{n conta registros que mencionam o tema em qualquer uma das cinco perguntas; categorias não exclusivas. Frequências completas, definições e exemplos estão no material complementar.}
\end{table*}
